\documentclass[10pt]{article}
\usepackage[T1]{fontenc}
\usepackage[utf8]{inputenc}
\usepackage[french]{babel}
\usepackage{amsmath,amssymb,mathrsfs,mathpazo}
\usepackage{geometry}
\geometry{a4paper,margin=2cm}
\usepackage{graphicx,array,tabularx,tabularray,colortbl,multirow,multicol,booktabs,enumitem,pifont,ulem,fancybox,diagbox,cancel}
\usepackage{pstricks,pstricks-add,pst-plot,pst-tree,pst-node,pst-text,pst-func,pst-math,pst-bspline,pst-3dplot,pst-eucl,pst-solides3d}
\let\pscalculate\undefined
\usepackage{pst-bezier}
\usepackage{tikz,pgfplots,tkz-tab}
\usetikzlibrary{arrows,arrows.meta,patterns,calc,positioning,trees,decorations.markings,shapes.geometric,shapes.misc,shapes.arrows,intersections,decorations.pathmorphing,decorations.pathreplacing,backgrounds,shadows,babel}
\pgfplotsset{compat=1.18}
\usepgfplotslibrary{fillbetween}
\usepackage[np]{numprint}
\usepackage{listings,makecell,pas-tableur,twoopt}
\let\barre\undefined
\usepackage[pstricks]{tablvar}
\definecolor{darkgreen}{RGB}{0,100,0}
\providecommand{\red}{\color{red}}
\providecommand{\blue}{\color{blue}}
\providecommand{\green}{\color{green}}
\providecommand{\black}{\color{black}}
\providecommand{\white}{\color{white}}
\providecommand{\gray}{\color{gray}}
\providecommand{\cyan}{\color{cyan}}
\providecommand{\magenta}{\color{magenta}}
\providecommand{\yellow}{\color{yellow}}
\providecommand{\violet}{\color{violet}}
\providecommand{\orange}{\color{orange}}
\providecommand{\euro}{\texteuro}
\providecommand{\EUR}{\texteuro}
\providecommand{\ieme}{\textsuperscript{e}}
\providecommand{\degres}{\ensuremath{^\circ}}
\providecommand{\degree}{\ensuremath{^\circ}}
\providecommand{\vv}[1]{\overrightarrow{#1}}
\providecommand{\Coord}[1]{#1}
\providecommand{\decosix}{}
\providecommand{\starredbullet}{\ensuremath{\bullet}}
\pagestyle{empty}
\setlength{\parindent}{0pt}
\usepackage{adjustbox}\DeclareUnicodeCharacter{2236}{:}\DeclareUnicodeCharacter{2190}{\ensuremath{\leftarrow}}
\providecommand{\R}{\mathbb{R}}
\providecommand{\N}{\mathbb{N}}
\providecommand{\Z}{\mathbb{Z}}
\providecommand{\Q}{\mathbb{Q}}
\providecommand{\D}{\mathbb{D}}
\providecommand{\C}{\mathbb{C}}
\providecommand{\e}{{\mathrm{e}}}
\providecommand{\vect}[1]{\overrightarrow{#1}}
\providecommand{\vectt}[1]{\overrightarrow{\mathrm{#1}}}
\providecommand{\vvt}[1]{\overrightarrow{\mathrm{#1}}}
\providecommand{\np}[2][]{{#2~#1}}
\providecommand{\vv}[1]{{\overrightarrow{#1}}}
\renewcommand{\d}{{\mathrm{d}}}
\providecommand{\degres}{{^\circ}}
\providecommand{\degre}{{^\circ}}
\providecommand{\blue}{\color{blue}}
\providecommand{\red}{\color{red}}
\providecommand{\enskip}{\quad}
\providecommand{\llbracket}{\lbrack\!\lbrack}
\providecommand{\rrbracket}{\rbrack\!\rbrack}
\providecommand{\euro}{{\text{€}}}
\providecommand{\no}{{\text{n°}}}
\providecommand{\interval}[2]{{[#1;#2]}}
\providecommand{\intervFO}[2]{{[#1;#2[}}
\providecommand{\dotfill}{\ldots}

\providecommand{\ds}{\displaystyle}
\providecommand{\textding}[1]{#1}
\providecommand{\Oij}{\(\mathrm{O};\vec{\imath},\vec{\jmath}\)}
\providecommand{\Oijk}{\(\mathrm{O};\vec{\imath},\vec{\jmath},\vec{k}\)}
\providecommand{\Ouv}{\(\mathrm{O};\vec{u},\vec{v}\)}
\providecommand{\intervFF}[2]{\ensuremath{[#1;#2]}}
\renewcommand{\euro}{\ifmmode\text{\texteuro}\else\texteuro\fi}
\DeclareUnicodeCharacter{00B0}{\ifmmode^{\circ}\else\textdegree\fi}
\DeclareUnicodeCharacter{20AC}{\ifmmode\text{\texteuro}\else\texteuro\fi}
% Énoncés : Éducation nationale. Transcription et corrigés : APMEP, Denis Vergès et contributeurs. Découpage et classement : CoopMaths, Sébastien Lozano, Éric Elter et contributeurs. Adaptation de la présentation : Maths Explorer.
% CC BY-SA : https://coopmaths.fr/www/about/
% https://coopmaths.fr/alea/static/bac/2021/tex/bac_2021_06_metropole_10_cor.tex
\begin{document}
\noindent\textit{APMEP / CoopMaths. CC BY-SA. Présentation adaptée par Maths Explorer.}\par\medskip
\smallskip

\begin{tabular}{|l|}\hline
Principaux domaines abordés :\\
Fonction exponentielle; dérivation; convexité\\ \hline
\end{tabular}

\begin{center}\textbf{Partie 1}\end{center}








\begin{enumerate}
\item
\begin{list}{\textbullet}{D'après la courbe représentant la fonction dérivée $f'$:}
\item la fonction $f'$ est positive sur $]-\infty~;~1[$ donc la fonction $f$ est croissante sur cet intervalle;
\item la fonction $f'$ est négative sur $]1~;~+\infty[$ donc la fonction $f$ est décroissante sur cet intervalle.
\end{list}

\item 
\begin{list}{\textbullet}{D'après la courbe représentant la fonction dérivée $f'$:}
\item la fonction $f'$ est décroissante sur $]-\infty~;~0[$ donc la fonction $f$ est concave sur cet intervalle; 
\item la fonction $f'$ est croissante sur $]0~;~+\infty[$ donc la fonction $f$ est convexe sur cet intervalle.
\end{list}
\end{enumerate}












\medskip

\begin{center}\textbf{Partie 2}\end{center}

\smallskip

On admet que la fonction $f$ mentionnée dans la Partie 1 est définie sur $\R$ par :
$f(x) = (x + 2)\e^{-x}.$





\medskip

\begin{enumerate}
\item Pour tout nombre réel $x$, $f(x) = (x + 2)\e^{-x}=x\e^{-x}+2\e^{-x}
= \dfrac{x}{\e^{x}}+ 2\e^{-x}$.



D'après le cours: $\displaystyle\lim_{x \to ++ \infty} \dfrac{\e^{x}}{x}= +\infty$ donc  $\displaystyle\lim_{x \to + \infty} \dfrac{x}{\e^{x}}= 0$.

De plus $\displaystyle\lim_{x \to + \infty} \e^{-x}= 0$ donc $\displaystyle\lim_{x \to + \infty} f(x) = 0$.



On en déduit que la courbe $\mathcal{C}$ admet la droite d'équation $y=0$, c'est-à-dire l'axe des abscisses, comme asymptote horizontale en $+\infty$.

On admet que $\displaystyle\lim_{x \to - \infty} f(x) = - \infty$.
\item 
	\begin{enumerate}
		\item
$f'(x)=1\times \e^{-x} + (x+2)\times (-1)\e^{-x} = (1-x-2) \e^{-x}=(-x-1)\e^{-x}$. 		
		
		\item
Pour tout $x$, $\e^{-x}>0$ donc $f'(x)$ est du signe de $-x-1$; donc $f'(x)$ s'annule et change de signe en $x=-1$.

$f(-1) = (-1+2)\e^{1}=\e$; on établit le tableau de variations de $f$ sur $\R$:

\begin{center}
{\renewcommand{\arraystretch}{1.5}
\psset{nodesep=3pt,arrowsize=2pt 3}
\def\esp{\hspace*{2.5cm}}
\def\hauteur{0pt}
$\begin{array}{|c|l*4{c}|}
\hline
x & -\infty  & \esp & -1 & \esp & +\infty \\ 
\hline
-x-1 &    &   \pmb{+} & \vline\hspace{-2.7pt}0 & \pmb{-} & \\ 
\hline
f'(x) &    &   \pmb{+} & \vline\hspace{-2.7pt}0 & \pmb{-} & \\ 
\hline
 &  &  &   \Rnode{max}{\e}  &  &  \rule{0pt}{\hauteur} \\  
f(x) &   &     &  &  &  \rule{0pt}{\hauteur} \\ 
 &  \Rnode{min1}{-\infty} &   &  &  &   \Rnode{min2}{0} \rule{0pt}{\hauteur}    
 \ncline{->}{min1}{max} 
 \ncline{->}{max}{min2} 
 \\ 
\hline
\end{array} $
}
\end{center}	
	
		\item
Sur l'intervalle $[-2~;~-1]$, la fonction $f$ est strictement croissante et continue car dérivable sur cetintervalle. $f(-2)=0<2$ et $f(-1)=\e>2$	donc, d'après le corollaire du théorème des valeurs intermédiaires, l'équation $f(x)=2$ admet une solution unique sur l'intervalle $[-2~;~-1]$.
		
		
	\end{enumerate}
\item
$f''(x)=(-1)\times \e^{-x} + (-x-1)\times (-1)\e^{-x} = (-1+x+1) \e^{-x}=x\e^{-x}$

$\e^{-x}>0$ pour tout $x$, donc $f''(x)$ est du signe de $x$.

\begin{list}{\textbullet}{}
\item Sur $]-\infty~;~0[$, $f''(x)<0$ donc la fonction $f$ est concave.
\item Sur $]0~;~+\infty[$, $f''(x)>0$ donc la fonction $f$ est convexe.
\item En $x=0$, la dérivée seconde s'annule et change de signe donc le point A d'abscisse 0 de $\mathcal{C}$ est le point d'inflexion de cette courbe.
\end{list}



\end{enumerate}

\end{document}
